Este capítulo cierra el trabajo frente a los objetivos del Capítulo~\ref{objetivos} y a las cifras de
evaluación leftover sobre NYIT. El detalle experimental, las figuras y las tablas están en el
Capítulo~\ref{chapter:06}; las limitaciones y las líneas abiertas, en el Capítulo~\ref{chapter:08}.

\section{Objetivo general}

Se ha diseñado, implementado y evaluado un verificador \textit{open-set} de autenticación continua y pasiva
en un smartphone público (dump NYIT \textit{home+lab cleaned}, iPhone XR), con fusión de dinámica de tecleo,
gesto táctil (\textit{swipe}) y gestos de desplazamiento. El sistema opera sobre ventanas de interacción
sin exigir un gesto de login y sin persistir el texto tecleado.

Ese objetivo se cumple como \textbf{demostrador de arquitectura y protocolo}, no como sustituto de un factor
biométrico de alta seguridad. El mejor sistema propuesto ---Transformer de tres modalidades--- alcanza AUC
$0{,}83$ en leftover de evaluación, de modo que el ranking genuino/impostor es viable. El EER $24{,}60\%$
no admite un umbral de baja fricción: en el punto de error igual se acepta y se rechaza del orden de uno
de cada cuatro intentos. Las estimaciones iniciales de EER $1$--$3\%$ (Capítulo~\ref{chap:desarrollo})
\textbf{no se materializan} en este corpus (tecleo solo \textit{press}, \textit{swipe} que no es rúbrica, fusión no
simultánea a nivel de ensayo).

\section{Objetivos específicos}

\begin{enumerate}
    \item \textbf{Arquitectura multimodal.} Un Transformer con encoders por modalidad y atención cruzada,
    entrenado con \textit{Set2set Loss}, produce la mejor representación conjunta del dump. El \textit{swipe}
    no aporta identidad en unimodal (EER $\approx 50\%$); el scroll sí reduce el EER del
    Transformer ($28{,}14\%\to 24{,}60\%$) y no ayuda al ensemble en leftover de tecleo.

    \item \textbf{Dos enfoques bajo el mismo protocolo.} Transformer y ensemble (BiLSTM / CNN-LSTM / MLP
    con fusión de \textit{scores}) se evalúan con los mismos splits \textit{user-disjoint} (49/16/18,
    \textit{seed} 42), $k=5$, leftover e impostor aleatorio. Sobre el leftover emparejado (545 pares),
    la atención cruzada (EER $24{,}60\%$) supera al ensemble meta ($27{,}72\%$) y al mejor RF ($26{,}78\%$).
    La diferencia con Cariello et al.\ no se explica por ``usar RF en vez de una red''.

    \item \textbf{Pérdidas de verificación.} \textit{Triplet Loss} en los expertos y \textit{Set2set} en
    el Transformer permiten comparar un perfil de enrolamiento contra consultas continuas. El umbral se
    fija en desarrollo (EER) y se aplica en evaluación.

    \item \textbf{Evaluación y referentes.} El corpus ejecutado es NYIT, no Aalto, HuMIdb ni BehavePassDB
    (estas últimas, no disponibles aquí; se dejan como trabajo futuro). TypeNet, Type2Branch y BehaveFormer
    se citan como SOTA \textbf{ajena}. El único referente del mismo set es Cariello et al.\ (RF/$k$-NN/SVM + mocap Vicon, otro protocolo). No hay un Transformer SOTA
    publicado en NYIT con el que igualar cifras.

    \item \textbf{Longitud de secuencia, detección y TTD.} En el BiLSTM de tecleo, $N\in\{25,50,100\}$ apenas mueve
    el EER ($48{,}88\%$, $47{,}56\%$, $46{,}48\%$): no hay salto tipo TypeNet. La tasa de detección ($k=2$ ventanas)
    es $97{,}7\%$ en el Transformer frente a $41{,}2\%$ en el ensemble; el TTD medio en detectados es $11{,}6$ y
    $24{,}2$ ventanas, respectivamente. No se modela tiempo de reloj de una sesión real.

    \item \textbf{Privacidad.} No se persistió texto de las pulsaciones. Las trazas versionables son
    métricas, incrustaciones y puntuaciones con identificadores seudónimos, alineado con minimización
    (RGPD) y \textit{Privacy by Design}.
\end{enumerate}

La IMU se documentó como experimento adicional (fusión tecleo+IMU: EER $34{,}63\%$). No forma parte del
sistema propuesto ni de su tabla de resultados.

\section{Conclusión}

En NYIT leftover emparejado (545 pares), el Transformer 3-mod es el mejor sistema propuesto; le siguen
el RF de rasgos 3-mod y el ensemble meta. Open-set continuo es viable en \textbf{ranking} y no lo es como login
operativo tipo Aalto o tipo Vicon. El trabajo aporta un pipeline reproducible, una comparación honesta
de arquitecturas (incluida una cabeza clásica en el mismo protocolo) y una lectura explícita de los
límites del dump. Lo que queda abierto ---DT real, rúbrica alineada al tecleo, impostor hábil, más dispositivos,
Aalto/HuMIdb/BehavePassDB con el mismo código--- se detalla en el capítulo siguiente.
